\section{Gestores de bases de datos}

Los gestores de bases de datos, también conocidos como sistemas de gestión de bases de datos (Database Management Systems, DBMS), son herramientas diseñadas para facilitar la administración, almacenamiento, consulta y manipulación de datos de manera estructurada \parencite{DBMSGFG}. Estos sistemas actúan como intermediarios entre los usuarios o aplicaciones y los motores de bases de datos, permitiendo una interacción más eficiente y controlada con la información (Ver figura \ref{fig:DBMS}).

\begin{figure}[H]
    \centering
    \includegraphics[width=0.5\linewidth]{teorico/dbms.png}
    \caption{Flujo de trabajo de un DBMS.}
    \label{fig:DBMS}
\end{figure}

Su principal objetivo es garantizar la persistencia de los datos a lo largo del tiempo, asegurando características fundamentales como la durabilidad, la integridad y la seguridad de la información. Para ello, los gestores de bases de datos incorporan mecanismos de control de acceso, respaldo, recuperación ante fallos y validación de datos.

Además, estos sistemas permiten gestionar entornos tanto mono-usuario como multiusuario, facilitando el acceso concurrente a la información sin comprometer su consistencia. También ofrecen opciones de configuración y personalización que permiten adaptarlos a las necesidades específicas de cada sistema, favoreciendo su integración dentro de distintas arquitecturas de software.

\subsection{PostgreSQL}

PostgreSQL es un sistema de gestión de bases de datos relacional de código abierto, reconocido por su robustez, extensibilidad y alto rendimiento en el manejo de grandes volúmenes de información. Se ha consolidado como una de las soluciones más utilizadas en la actualidad debido a su capacidad para ejecutar consultas complejas y gestionar datos de manera eficiente y confiable \parencite{PostgreSQL}.

Este gestor se caracteriza por contar con un motor potente que garantiza la integridad, consistencia y seguridad de los datos, además de ofrecer soporte para múltiples tipos de datos, funciones avanzadas y mecanismos de optimización de consultas. Gracias a estas características, PostgreSQL ha sido adoptado por numerosas organizaciones y empresas para el almacenamiento y gestión de información crítica.

PostgreSQL utiliza el lenguaje SQL (Structured Query Language) como medio principal para la definición, manipulación y consulta de datos. A través de este lenguaje, es posible realizar operaciones como inserciones, actualizaciones, eliminaciones y consultas sobre tablas y relaciones, siguiendo el modelo de bases de datos relacionales. Esto permite manejar grandes volúmenes de datos bajo estructuras organizadas, facilitando su acceso y administración de manera eficiente.